\scanpage{74}
% Archive pencil foliation: 73.

\handwritten{7.} Nous verrons çi-dessus que la réponse à la question précédente est affirmative lorsque $k$ est un corps fini, en particulier si c'est le c\unsure{\deleted r o|\deleted o r}\unsure{\deleted o r|\deleted r o}ps premier $\underline{F}_p$, On trouve alors, grâce à cette structure d'extension sur $\underline{G}(\underline{F}_p)$, un homomorphisme canonique de Gerbes (pour tout $k$)
\[
\underline{G}(k)\handwritten{\simeq}\underline{G}(\underline{F}_p)\mathrm{x}_{\deleted{\unsure{\underline{G}(}}\handwritten{\unsure{\underline{\pi}|\underline{\Pi}}}_p(\underline{F}_p))}\handwritten{\underline{\pi}}_p(k)\ ,
\]
et pour $k$ parfait, une réponse affirmative à la question précédente revient à l'assertion que l'homomorphisme précédent est une équivalence de Gerbes. Comme on connait parfaitement bien la Gerbe $\underline{G}(\underline{F}_p)$ et son homomorphisme dans $\handwritten{\underline{\pi}}_p(\underline{F}_p)$ (cf plus bas) cela donnerait donc la structure de $\underline{G}(k)$ en termes de la connai\unsure{\deleted a s|\deleted s a}sance de \deleted{${}_p(\underline{F}_p$} $\handwritten{\underline{\unsure{\Pi|\pi}}}_{\deleted{\handwritten{\unsure{\alpha|a|\ell}}}p}\unsure{\deleted{(}(|(}k)$ (connu quand\handwritten{/}on connait $\handwritten{\pi}_1(k)=\operatorname{Gal}(\handwritten{\bar{k}}/k)$) et \deleted{de} l'homomorphisme $\handwritten{\underline{\pi}}_p(k\deleted{\unsure{F_p|Fr}})\handwritten{\longrightarrow\underline{\pi}}_p(\underline{F}_p)$ (connu quand on connait $\handwritten{\pi}_1(k)\handwritten{\longrightarrow\pi}_1(\underline{\unsure{\deleted K F|\mathbb{F}}}_p)\handwritten{\longrightarrow}\handwritten{\hat{\underline{Z}}}$). Cela résoudrait donc la question de la structure de $\underline{G}(k)$ pour $k$ parfait. On notera que l'objet de Tate de $\mathrm{Fc}\unsure{\deleted{(}\mathrm r|\mathrm r}\mathrm{iso}(k)$ provient de celui \unsure{\deleted a d|\deleted d a}e $\mathrm{Fcriso}(\underline{F}_p)$ et sera donc également déterminé. Le foncteur fibre canonique sur $K$\inserted{$=K(k)$} correspond alors à un foncteur fibre de $\mathrm{Fcriso}(\underline{F}_p)$ sur $K$, un foncteur fibre de \deleted{$\underline{\unsure{\pi|\Pi}}_p(k)$}\inserted{la catégori\unsure{\deleted t e|\.{e}|e} des $\underline{Q}_p$-faisceaux constr. s\unsure{\deleted t u|\deleted u u}r $k$} sur $K$, et un isomorphisme entre les ``restrictions'' de ces foncteurs à la catégorie des $\underline{Q}_p$-faisceaux constructibles sur $\underline{F}_p$ ; or le premier foncteur fibre n'\unsure{\deleted t e|\.{e}|e}\unsure{\deleted d s.|\deleted d s|s.|s}t autre que \inserted{celui déduit par} l'\unsure{\deleted t e|\.{e}|e}xtension des scalaire\unsure{\deleted a s|\deleted s a|\deleted s s} $\underline{Q}_p=K\handwritten{\unsure{\deleted{-}(|(}}\underline{F}_p)\handwritten{\longrightarrow}K=K(k)$ du foncteur fibre analogue sur $\unsure{\deleted{\mathrm{G}}\mathrm{F}|\mathbb{F}|\mathrm{F}|\mathrm{G}}\mathrm{criso}(\underline{F}_p)$, le deuxième est \unsure{\deleted n l|\deleted h l|l|h}e foncteur \inserted{\handwritten{$i$:}} $V\handwritten{\mapsto}V\otimes_{\underline{Q}_p}K$, et l'isomorphisme \deleted{canonique} entre les restrictions est l'isomorphisme canonique évident\inserted{\handwritten{(compatibilité de $i$ \unsure{aux|avec} \unsure{changements|changement} de corps de base)}}. Donc la structure supplémentaire de $\underline{G}(k)$ (savoir sa section \deleted{sur} canonique sur $K$) est déterminée par la connaissance de la section de $\handwritten{\underline{\pi}}_p(k)$ sur $K$ (i.e. le foncteur précédent $V\handwritten{\mapsto}V\otimes_{\underline{Q}_p}K$ ) et l'isomorphisme de l'i\unsure{\deleted a m|\deleted m a}\unsure{\deleted g a|\deleted a g}\unsure{\deleted m g|\deleted g m}e de ladite section dans $\handwritten{\underline{\pi}}_p(\underline{F}_p)$ avec la section provenant de\deleted{\unsure{l|s|a}} la section canonique de $\handwritten{\underline{\pi}}_p(\underline{F}_p)$ sur $\underline{Q}_p$ ; ce sont là des données qui ne sont évidemment pas contenues dans celle du groupe

\margin{\handwritten{\uline{NB} On a donc, en comparant le foncteur ``fibre ``\deleted{\unsure{oubli|à la|à k}} en $\unsure{\bar{k}|k}$'' sur $Q_p$, et \unsure{le|les} \unsure{foncteur|foncteurs} \unsure{précédent|précédents} un torseur canonique sur $K$ de groupe \deleted{\unsure{de|et}}(lien $\underline{\pi}_p(k))_K$\deleted{\unsure{relatif à|relatif sur} $\unsure{Q_p|\mathbb{Q}_p}$}}}
